\section{Introduction}

Neural network weights encode architectural priors and training history. A ResNet-50's 25 million parameters capture BatchNorm statistics, residual connection patterns, and learned kernel structures that distinguish it from a ViT-Base with patch embeddings and attention mechanisms. Can we predict a model's test accuracy or detect backdoor tampering just from weight tensors—without executing the model? Weight-space learning—treating model parameters as first-class data \cite{schurholt2022hyper}—enables meta-learning applications like model zoo search, neural architecture prediction, and anomaly detection.

Processing weights as data requires handling variable layer shapes and high dimensionality while preserving architectural signal. A standard approach flattens each layer into tokens for sequence models, but systematic validation of this tokenization strategy is missing. Does layer-wise processing retain sufficient structure for property prediction, or do cross-layer dependencies dominate? Architecture family classification serves as a foundational validation: if tokenization loses critical structure, even coarse-grained family discrimination (ResNet vs ViT) should fail.

Prior work in weight-space learning demonstrates feasibility through task performance—hyper-representations transfer across meta-learning tasks \cite{schurholt2022hyper}, model stitching shows layer-wise analysis validity \cite{lenc2015understanding}, and weight statistics correlate with model properties \cite{unterthiner2020predicting}. These approaches establish that weight-space features carry signal but lack controlled experiments comparing tokenization strategies (statistical aggregation vs learned token embeddings) independent of downstream task complexity. The foundational question remains: \textit{How do different layer-wise processing approaches compare in preserving structural signal for property prediction?}

We address this gap through systematic validation on timm Model Zoo (100 pretrained vision models across 4 architecture families: ResNet, ViT, EfficientNet, ConvNeXt). Our contributions:

\begin{enumerate}
\item \textbf{Tokenization Strategy Validation}: Layer-wise processing (flatten + pad + per-layer normalize) achieves 80\% $\pm$ 4.2\% accuracy on architecture family classification, significantly exceeding our 60\% gate threshold and 25\% random baseline.

\item \textbf{Baseline Comparison}: Simple per-layer statistical aggregation (mean/std/L2 norm) matches transformer-based token embedding performance (both 80\% $\pm$ 4.2\%) on small datasets, revealing that family-level discriminability doesn't require cross-layer attention modeling when architecture families exhibit strong layer-level separability.

\item \textbf{Practical Method}: We provide a validated tokenization pattern—flatten weight matrices, zero-pad to fixed length, apply per-layer normalization—tested on real pretrained model zoo with reproducible hyperparameters and evaluation protocol.
\end{enumerate}

Our findings validate layer-wise weight processing as a viable foundation for weight-space transformers while highlighting that simple statistical baselines provide strong comparison points. This work enables future research on larger datasets and harder downstream tasks (backdoor detection, property prediction, model synthesis) with confidence that layer-wise processing approaches preserve necessary structural signal.
